\section{A linear regularization that erases its endpoint constant}\label{sec:mono}
Arithmetic holonomy suggests studying a rational germ that gains
arithmetic information if its endpoint period is rational; see the
framework described by Calegari, Dimitrov and Tang~\cite{CDT}.
No holonomy bound is applied here. The following obstruction is an
elementary monodromy statement about one specified primitive.

Let $E=x\,\mathrm d/\mathrm dx$ and
$F(x)=\sum_{m\geq1}\chi_4(m)x^m/m^4$. Put $F_j=E^jF$ and
\begin{equation}\label{eq:regularized}
 C(x)=F_0-\log x\,F_1+\frac{\log^2x}{2}F_2
                              -\frac{\log^3x}{6}F_3.
\end{equation}
Since $E^4F=x/(1+x^2)$, differentiation telescopes:
\[
 C'(x)=-\frac{\log^3x}{6(1+x^2)},\quad
 C(x)=\int_0^x\frac{-\log^3t}{6(1+t^2)}\dd t,\quad C(1)=\theta.
\]
For any rational $q$, $C_q=C+q-\theta$ has a rational Taylor germ at
$x=1$, including its constant term. This is a full rational-jet repair,
not just a rational value. It does not remove the logarithmic branch
at zero.

\begin{proposition}[Scoped monodromy obstruction; E04]\label{prop:mono}
Let $L$ be a linear differential operator with rational-function
coefficients in $x$, even over $\C$. If $L(C_q)$ is meromorphic near
zero, then $L(1)=0$. In particular $L(C_q)=L(C)$ is independent of
the additive constant $q-\theta$.
\end{proposition}
\begin{proof}
Let $M_0$ denote continuation once positively about zero and
$\Delta=(M_0-1)C_q$. All $F_j$ are holomorphic at zero, and
\[
 \Delta'(x)=-\frac{(\log x+2\pi i)^3-\log^3x}{6(1+x^2)}.
\]
Meromorphicity implies $L(\Delta)=0$. Continue to $x=i$ on the
branch with $\log i=\pi i/2$. Then
\[
 \Res_{x=i}\Delta'
 =-\frac{(5\pi i/2)^3-(\pi i/2)^3}{12i}
 =\frac{31\pi^3}{24}\ne0.
\]
A small circuit about $i$ adds $31i\pi^4/12$ to $\Delta$.
The coefficients of $L$ are single-valued. Continuing $L(\Delta)=0$
and subtracting thus gives $(31i\pi^4/12)L(1)=0$.
Paths can avoid the finite set of poles of $L$.
\end{proof}

The symmetry $C(1/x)=C(x)$ follows by differentiation and equality at
$x=1$. With $x=(1-z)/(1+z)$ and $y=z^2$ it yields
\[
 C(x)=\theta+J(y),\qquad
 J'(y)=-\frac23\frac{yA(y)^3}{1+y},\qquad
 A(y)=\frac{\atanh\sqrt y}{\sqrt y},\quad J(0)=0.
\]
Writing $B=yA$ gives the rational system
\[
 (B^k)'=\frac{kB^k}{2y}+\frac{kB^{k-1}}{2(1-y)}\ (k=1,2,3),
 \qquad J'=-\frac{2B^3}{3y^2(1+y)}.
\]
The span of $1,B,B^2,B^3,J$ is unchanged by replacing $J$ with
$q+J$. It has dimension five: $B$ has infinitely many monodromy
values about $1$, so is transcendental over $\C(y)$, while $J'$
has nonzero residue $\pi^3/96$ at $y=-1$, where the first four
functions are locally meromorphic. Adding a rational constant changes
neither this module nor its continuation domains. These observations
explain this proposed shortcut's failure; they do not exclude
nonlinear expressions, extra functions or different differential systems.
